% chapters/36_openclaw_os.tex
% Part of: AI / LLM / Agents / Hardware Notes

\section{OpenClaw \& Agent OS Integration}

\subsection{OpenClaw architecture overview}

The NVIDIA stack for enterprise agentic deployment is composed of four distinct layers,
each with a different responsibility:

\begin{itemize}
  \item \textbf{OpenClaw} --- the agent/application layer. This is the thing that plans,
    calls tools, writes code, reads files, and tries to complete tasks. Think of it as
    the agent ``brain'' and behaviour loop.

  \item \textbf{OpenShell} --- the runtime governor. It sits \emph{between} OpenClaw and
    the real machine, network, and models, and decides what the agent is actually allowed
    to do. NVIDIA describes it as governing how the agent executes, what it can see and
    do, and where inference is routed.

  \item \textbf{NemoClaw} --- the packaged enterprise distribution. It takes OpenClaw and
    adds the missing operational pieces: safer setup, policy controls, privacy routing,
    model plumbing, and sandboxing. Installed in a single command via NVIDIA Agent
    Toolkit software.

  \item \textbf{Nemotron} --- NVIDIA's family of open models (reasoning, coding, visual
    understanding, safety, and agentic workflows). In this stack, Nemotron is the
    local or NVIDIA-hosted model option that NemoClaw/OpenShell can route inference to,
    especially when enterprises want on-prem execution.
\end{itemize}

A useful one-line shorthand for each layer:
\begin{center}
  OpenClaw = agent behaviour \quad
  Nemotron = model family \quad
  OpenShell = execution/security runtime \quad
  NemoClaw = packaged deployment stack
\end{center}

The key design philosophy is that ``wrapper'' is the right mental model for NemoClaw:
it does not replace the agent brain; it adds a controlled execution environment
\emph{around} it. The core question OpenClaw answers is \emph{``how do we make agents
capable?''}; OpenShell answers \emph{``how do we make capable agents safe to run?''}
NemoClaw bundles both answers into something an enterprise can actually operate.

\subsection{OS-level agent primitives}

OpenShell exposes three core architectural blocks that together form an agent OS:

\begin{enumerate}
  \item \textbf{Sandbox} --- the isolated environment where the agent runs. Mistakes or
    malicious/confused behaviour are contained inside the sandbox boundary and cannot
    reach the host OS or other sessions.

  \item \textbf{Policy engine} --- checks what actions are allowed across the filesystem,
    network, and process layers before any action executes. Rules are defined by the
    operator/enterprise, not the agent.

  \item \textbf{Privacy router} --- decides whether a given inference request or data
    payload should stay on a local model (e.g.\ a local Nemotron instance) or may be
    forwarded to a frontier cloud model. The routing decision is based on operator
    cost/privacy policy, explicitly \emph{not} on what the agent wants to do.
\end{enumerate}

\textbf{Action interception flow.} When OpenClaw attempts any action (read a file, call
a tool, invoke a model), OpenShell intercepts it before execution:
\begin{enumerate}
  \item Agent requests an action.
  \item Policy engine evaluates whether the action is permitted under current rules.
  \item Privacy router determines where any associated inference should go.
  \item Action executes only if policy allows it; otherwise it is blocked and the
    agent receives an error or a constrained alternative.
\end{enumerate}

\textbf{Example policy rules:}
\begin{itemize}
  \item HR data must stay on-device; only an approved local Nemotron model may see it.
  \item Code repository content may be routed to an approved cloud model.
  \item Outbound network requests are limited to a specified domain whitelist.
  \item The agent may read \texttt{/workspace/invoices} but not \texttt{/home/user/.ssh}.
  \item Only certain folder paths are writable; cloud model calls require scrubbed input.
\end{itemize}

\subsection{Filesystem \& process access}

OpenShell implements controls at the OS primitive level, not just through soft system
prompt instructions. Four categories of guardrails are enforced at runtime:

\begin{itemize}
  \item \textbf{Filesystem limits.} Scoped read/write access per session. An agent
    handling support tickets can be restricted to \texttt{/tickets/current} and
    explicitly denied access to credential stores or other user home directories.

  \item \textbf{Blocked network egress.} Outbound connections require policy approval.
    An enterprise can prevent customer PII from leaving the perimeter by blocking
    unapproved external endpoints at the network namespace level.

  \item \textbf{Process restrictions.} Limits on what processes the agent can spawn
    (e.g.\ no arbitrary shell exec, no privilege escalation).

  \item \textbf{Controlled inference backends.} The agent cannot freely choose which
    model API to call; the privacy router intercepts model calls and enforces
    routing policy.
\end{itemize}

Specific kernel-level mechanisms cited in NVIDIA's OpenShell documentation:
\textbf{Landlock} (Linux kernel filesystem access control), \textbf{seccomp}
(system call filtering), and \textbf{network namespaces} (network isolation per
process group). These are the same primitives used in container runtimes like
Docker/gVisor, adapted for agentic workloads.

\subsection{Security sandboxing}

The most important design idea in OpenShell is \textbf{out-of-process policy
enforcement}: moving the real control point \emph{outside} the agent, so the agent
cannot override it even if its reasoning is compromised, confused, or
prompt-injected.

\textbf{The problem with in-agent safety.} Most agent safety today is ``inside the
agent'': system prompts, behavioural rules, tool descriptions that say ``don't do X''.
These soft protections can fail if the agent loses context, is prompt-injected via
observed content, or simply reasons its way around them. In-context instructions
cannot enforce themselves.

\textbf{OpenShell's answer.} By intercepting actions at the OS/runtime layer before
they execute, OpenShell provides hard enforcement that the agent model cannot
circumvent. NVIDIA compares this to a browser security model: sessions are isolated,
and permissions are checked at the execution boundary, not just in the agent's
``mind''. The agent may \emph{believe} an action is permitted; OpenShell decides
whether it actually executes.

This is why OpenShell is better described as:
\begin{itemize}
  \item a runtime governor
  \item a policy firewall
  \item a sandbox manager
  \item an inference router
\end{itemize}
\ldots rather than as a model, a chatbot, or a prompt-engineering layer.

\textbf{Kernel-level isolation.} NVIDIA explicitly targets kernel-level sandboxing
(Landlock, seccomp, network namespaces) rather than process-level or
container-level approaches, giving stronger isolation guarantees for long-running
autonomous sessions.

\subsection{Desktop automation patterns}

\textbf{Enterprise deployment example: bank support ticket agent.}

Without a controlled runtime, an unconstrained agent processing support tickets might:
\begin{itemize}
  \item read too many internal files beyond those needed for the task
  \item leak customer PII to a cloud model in its prompt
  \item execute unsafe shell commands with unexpected side-effects
  \item POST data to an unapproved external endpoint
\end{itemize}

With NemoClaw/OpenShell, the operator configures policy such that:
\begin{itemize}
  \item customer PII stays on-prem; only an approved local Nemotron model sees raw
    ticket text
  \item only summarised, scrubbed data may be forwarded to a frontier cloud model
  \item the agent may access \texttt{/tickets/current} but not engineer laptops or
    credential folders
  \item outbound network access is limited to an approved internal API whitelist
\end{itemize}

\textbf{The enterprise value proposition} is therefore not ``make the agent smarter''
but \textbf{``make the agent deployable''}. The big blocker for long-running agents in
production is not model quality --- it is \emph{trust}. NVIDIA frames this explicitly:
the missing primitives are sandboxing, permissions, and isolation. If a safe runtime
layer exists, enterprises can let agents run longer, touch more systems, and do useful
work without granting unconstrained access.

\textbf{Comparison to the Cowork/Claude Code approach.}
Both NemoClaw/OpenShell and Cowork follow the same fundamental pattern --- a VM or
sandbox separates the agent from the host, and policy tiers control what the agent can
reach --- but at different deployment targets:
\begin{itemize}
  \item Cowork targets individual/team desktop users: single-session Ubuntu VM,
    user-selected workspace mount, product-level prohibited-action lists.
  \item NemoClaw/OpenShell targets enterprise IT: kernel-level isolation, org-defined
    policy YAML, privacy routing between on-prem and cloud models, multi-session scale.
\end{itemize}

